\chapter{Glossário}
\label{cap:glossario}

Este apêndice reúne os termos técnicos empregados no livro, com a definição usada aqui e
o original em inglês. A lista de siglas encontra-se nas páginas pré-textuais.

\section{Operação de poços e escoamento}

\begin{description}[leftmargin=0pt,style=nextline,itemsep=0.35em]

\item[Anomalia] Comportamento do poço que se desvia do regime operacional esperado. Não
é sinônimo de falha: uma anomalia pode ser benigna, e uma falha pode ocorrer sem
anomalia detectável nos sensores disponíveis.

\item[Árvore de natal molhada (\ing{wet christmas tree})] Conjunto de válvulas instalado
na cabeça do poço submarino, que controla o fluxo entre o poço e as linhas de produção.

\item[Completação inteligente (\ing{intelligent completion})] Arranjo de poço com
válvulas de controle de intervalo operáveis remotamente, que permite regular a
contribuição de cada zona produtora sem intervenção física.

\item[Curva de formação de hidrato] Fronteira no plano pressão--temperatura que separa a
região onde hidratos de gás são termodinamicamente estáveis da região onde não são.

\item[Elevação por gás (\ing{gas lift})] Técnica de elevação artificial em que gás é
injetado na coluna de produção para reduzir a densidade da mistura e restabelecer o
fluxo.

\item[Golfada (\ing{slugging})] Regime de escoamento intermitente em que bolsões de gás
e de líquido alternam-se, produzindo oscilações de pressão e vazão.

\item[Hidrato de gás] Sólido cristalino formado por moléculas de gás aprisionadas em uma
estrutura de água, estável a alta pressão e baixa temperatura. Pode obstruir linhas de
produção.

\item[Incrustação (\ing{scaling})] Deposição progressiva de sais ou parafinas nas
paredes da tubulação ou na válvula reguladora, que restringe o fluxo ao longo de semanas
ou meses.

\item[Intervenção] Operação física sobre o poço, geralmente com embarcação dedicada, de
custo elevado. Reduzir intervenções não programadas é um dos objetivos do monitoramento.

\item[Linha de serviço] Linha secundária, distinta da linha de produção, usada para
injeção de químicos, gás de elevação ou operações auxiliares.

\item[Perda de produtividade] Redução sustentada da vazão produzida sob condições
operacionais equivalentes.

\item[Poço injetor] Poço destinado a injetar fluido no reservatório, para manutenção de
pressão ou varrido, em oposição ao poço produtor.

\item[Pré-sal] Província petrolífera brasileira situada abaixo de uma camada de sal, em
lâminas d'água tipicamente superiores a dois mil metros.

\item[Regime permanente (\ing{steady state})] Condição em que as variáveis do processo
permanecem aproximadamente constantes, em oposição ao transiente.

\item[Tempo não produtivo] Período em que o poço não produz por razão não programada.

\item[Transiente] Período de variação rápida das variáveis do processo, tipicamente após
uma manobra deliberada, durante o qual os modelos de regime permanente não se aplicam.

\item[Válvula de controle de intervalo] Válvula de fundo que regula a contribuição de
uma zona produtora específica.

\item[Válvula de segurança de subsuperfície] Válvula de fundo que fecha automaticamente
em condição de emergência, isolando o poço.

\item[Válvula reguladora de produção (\ing{production choke})] Válvula que controla a
vazão do poço para a linha de produção.

\end{description}

\section{Aprendizado de máquina}

\begin{description}[leftmargin=0pt,style=nextline,itemsep=0.35em]

\item[Acurácia balanceada] Média das revocações por classe. Usada quando as classes são
desbalanceadas, situação em que a acurácia simples é enganosa.

\item[Agrupamento (\ing{clustering})] Tarefa não supervisionada de particionar um
conjunto de dados em grupos de elementos semelhantes.

\item[Aprendizado de mundo aberto (\ing{open-world learning})] Paradigma em que o modelo
deve reconhecer que uma amostra não pertence a nenhuma classe conhecida, organizar as
amostras rejeitadas e incorporá-las como novas classes.

\item[Aprendizado de mundo fechado (\ing{closed-world learning})] Paradigma que pressupõe
que todas as classes que ocorrerão em operação estavam presentes no treinamento.

\item[Atenção (\ing{attention})] Mecanismo que pondera dinamicamente a contribuição de
cada posição de uma sequência, permitindo caminho de comprimento constante entre
posições distantes.

\item[Autoencoder] Rede treinada para reconstruir a própria entrada através de um
gargalo, aprendendo uma representação comprimida sem supervisão.

\item[Classe virtual] Agrupamento de amostras rejeitadas que satisfaz critérios de
coesão e é tratado provisoriamente como classe, até validação por especialista.

\item[Classe sintomática] Classe definida pelo sintoma observável e não pelo mecanismo
físico subjacente. Instâncias de uma mesma classe sintomática podem ter causas distintas,
o que compromete a coerência interna do rótulo.

\item[Deriva (\ing{drift})] Mudança, ao longo do tempo, da distribuição dos dados
(deriva de dados) ou da relação entre entrada e rótulo (deriva de conceito).

\item[Detecção de novidade (\ing{novelty detection})] Tarefa de identificar amostras que
não pertencem a nenhuma classe conhecida.

\item[Distância de Mahalanobis] Distância entre um ponto e uma distribuição, normalizada
pela covariância. É invariante a escala e sensível à correlação entre variáveis.

\item[Ensemble] Conjunto de modelos cujas saídas são combinadas. Neste livro, designa
sobretudo o conjunto de classificadores binários um-contra-todos.

\item[Época (\ing{epoch})] Uma passagem completa pelo conjunto de treinamento.

\item[Escore-$z$ (\ing{z-score})] Padronização que subtrai a média e divide pelo desvio
padrão.

\item[Espaço latente] Espaço de representação intermediária aprendido por um modelo. Sua
geometria determina o que métodos baseados em distância conseguem detectar.

\item[Índice de Calinski-Harabasz] Razão entre dispersão entre grupos e dispersão dentro
dos grupos. Valores maiores indicam melhor separação.

\item[Índice de Davies-Bouldin] Média da similaridade de cada grupo com seu vizinho mais
próximo. Valores menores indicam melhor separação.

\item[Interseção sobre união] Razão entre a interseção e a união de duas regiões. Usada
para avaliar a qualidade de segmentações.

\item[Janela deslizante (\ing{sliding window})] Recorte de comprimento fixo que avança
sobre a série temporal com passo definido.

\item[Mistura de especialistas] Arquitetura em que apenas um subconjunto dos parâmetros
é ativado por entrada, permitindo modelos de grande capacidade com custo de inferência
reduzido.

\item[Parada antecipada (\ing{early stopping})] Interrupção do treinamento quando a
métrica de validação deixa de melhorar por um número definido de épocas.

\item[Perceptron de múltiplas camadas] Rede neural composta por camadas densas
sucessivas com não linearidades.

\item[Rede recorrente] Rede que processa sequências mantendo um estado interno atualizado
a cada passo.

\item[Regularização por abandono (\ing{dropout})] Técnica que desativa aleatoriamente
unidades durante o treinamento, reduzindo sobreajuste.

\item[Reconhecimento em conjunto aberto (\ing{open-set recognition})] Formulação em que o
modelo pode recusar-se a atribuir uma classe, mas não organiza nem incorpora as amostras
recusadas.

\item[Risco de espaço aberto] Formalização do risco de classificar como conhecida uma
região do espaço de características arbitrariamente distante dos dados de treinamento.

\item[Segmentação de séries temporais] Tarefa de atribuir um rótulo a cada instante da
série, em oposição a atribuir um rótulo a uma janela inteira.

\item[Silhueta] Medida, por amostra, de quão bem ela se ajusta ao próprio grupo em
comparação ao grupo mais próximo. Varia em $[-1,1]$.

\item[Sobreajuste (\ing{overfitting})] Ajuste excessivo aos dados de treinamento, com
perda de desempenho em dados novos.

\item[Um-contra-todos (\ing{one-vs-all})] Estratégia que decompõe um problema
multiclasse em vários problemas binários independentes.

\item[Validação cruzada estratificada] Partição em dobras que preserva a proporção das
classes em cada dobra.

\end{description}

\section{Modelos de linguagem}

\begin{description}[leftmargin=0pt,style=nextline,itemsep=0.35em]

\item[Alucinação (\ing{hallucination})] Produção de afirmação fluente e plausível que
não é sustentada pela entrada nem por fato verificável.

\item[Cadeia de raciocínio (\ing{chain-of-thought})] Texto intermediário produzido antes
da resposta final. Não deve ser confundido com explicação do processo computacional
efetivo.

\item[Calibração] Grau de correspondência entre a confiança declarada por um modelo e a
frequência empírica de acerto.

\item[Contaminação de dados] Presença, no corpus de pré-treinamento, de dados usados na
avaliação, o que invalida a interpretação dos resultados como generalização.

\item[Instrução de sistema (\ing{system prompt})] Instrução que define papel, restrições
e formato de saída, aplicada antes da entrada específica.

\item[Prompt] Entrada textual fornecida ao modelo de linguagem, incluindo instrução,
contexto e dados.

\item[Saída estruturada] Resposta em formato especificado, tipicamente JSON, que pode ser
validada programaticamente.

\item[Temperatura] Parâmetro que controla a aleatoriedade da amostragem. Valores baixos
tornam a saída mais determinística.

\item[Transformador (\ing{transformer})] Arquitetura baseada em atenção, sem recorrência,
que permite paralelização do processamento da sequência.

\end{description}

\section{Métricas e avaliação}

\begin{description}[leftmargin=0pt,style=nextline,itemsep=0.35em]

\item[Detecção antecipada] Fração de eventos detectados antes do instante anotado pelo
especialista. Métrica da família SoftED.

\item[Erro percentual absoluto médio simétrico] Métrica de erro de reconstrução limitada
e simétrica em relação a sobre e subestimação.

\item[Especificidade] Fração de negativos corretamente identificados.

\item[Intervalo de Wilson] Intervalo de confiança para proporção binomial, adequado a
amostras pequenas e a proporções próximas de zero ou um.

\item[Medida $F_1$] Média harmônica entre precisão e revocação.

\item[Precisão] Fração de predições positivas que estão corretas.

\item[Revocação] Fração de positivos reais que foram detectados. Também chamada de
sensibilidade.

\item[SoftED] Família de métricas de detecção de eventos com tolerância temporal, que
credita parcialmente detecções próximas ao instante anotado.

\item[Top-$k$] Métrica que considera acerto se a resposta correta está entre as $k$
primeiras posições de uma lista ordenada.

\end{description}
